\begin{textsample}{POS Dim 2 – 2014\_c4 – Score 11.00 – 2014\_c4/enade\_2014\_letras\_por\_6f9b30ff67.txt}  \label{ex:f2_pos_120}
CAPÍTULO 4 PERCEPÇÃO DA PROVA As análises feitas neste capítulo tratam das percepções dos concluintes da Área de Letras - Português sobre a prova aplicada no ENADE/2014.
Estas percepções foram mensuradas por meio de nove questões que avaliaram desde o grau de dificuldade da prova até o tempo gasto para concluí -la.
As percepções sobre a prova foram relacionadas com o desempenho dos estudantes e com a Grande Região de funcionamento do curso.
O \textbf{questionário} de percepção da prova encontra -se ao final do Anexo IV, que traz a \textbf{reprodução} da prova.
O desempenho dos estudantes foi classificado em quatro quartos.
Para tanto, esse desempenho foi \textbf{ordenado} de forma ascendente.
O percentil 25, P25, também conhecido como primeiro quartil, é a nota de desempenho que deixa um quarto ( 25 \% ) dos valores observados abaixo e três quartos acima.
A Figura 1 apresenta uma ilustração deste conceito.
O quarto inferior de desempenho é composto pelas notas abaixo do primeiro quartil.
Já o percentil 75, P75, também conhecido como terceiro quartil, é o valor para o qual há três quartos ( 75 \% ) dos dados abaixo e um quarto acima dele.
O quarto superior de desempenho é composto pelas notas iguais ou acima do terceiro quartil.
O percentil 50, P50, também conhecido como mediana, é o valor que divide as notas em dois conjuntos de igual tamanho.
O segundo quarto inclui valores entre o primeiro quartil ( P25 ) e a mediana.
O terceiro quarto contém os valores entre a mediana ( P50 ) e o terceiro quartil ( P75 ).
Vale ressaltar que percentis, quartis e medianas são pontos que não obrigatoriamente pertencem ao conjunto original de dados, ao passo que os quartos são \textbf{subconjuntos} dos dados originais.
A seguir, serão apresentados gráficos com resultados selecionados, relativos às nove questões avaliadas por grupos de estudantes.
Os gráficos apresentam nas barras o percentual de alunos que assinalaram uma das opções ou a soma das porcentagens daqueles que assinalaram duas ( ou três ) delas.
Por exemplo, para as questões 1 e 2, os gráficos apresentam a porcentagem total de participantes que assinalaram as opções ( D ) difícil e ( E ) muito difícil.
Em cada barra foram assinalados também os extremos do intervalo de confiança de 95 \% como linhas verticais unidas por uma linha horizontal na forma da letra H maiúscula.
As Tabelas no Anexo II apresentam os valores absolutos e a distribuição percentual21 das alternativas válidas das nove questões, segundo o mesmo \textbf{recorte} de desempenho dos alunos e Grande Região de funcionamento do curso. 4.1 GRAU DE DIFICULDADE DA PROVA 4.1.1 Componente de Formação Geral Ao avaliarem “ Qual o grau de dificuldade desta prova na parte de Formação Geral? ” ( Questão 1 ), 33,0 \% do grupo de inscritos e presentes optaram pelas alternativas difícil ou muito difícil.
Entretanto, para mais da metade dos estudantes ( 56,8 \% ), o Componente de Formação Geral da prova foi considerado com grau de dificuldade médio ( Gráfico 4.1, Gráfico 4.2 e, no Anexo II, a Tabela II.1 ).
O percentual de estudantes que consideraram a prova como difícil ou muito difícil foi maior na região Sul, onde a proporção foi de 42,0 \%, enquanto a de menor incidência foi a Sudeste, com 25,7 \%.
No Gráfico 4.1, é possível observar que esta diferença é estatisticamente significativa.
Nas Grandes Regiões, a proporção de presentes à prova que consideraram o Componente de Formação Geral como sendo de grau de dificuldade médio esteve entre 51,0 \% na região Sul e 58,6 \% na região Nordeste. 21 Cumpre lembrar uma das convenções para tabelas \textbf{numéricas} ( pág. iii ) sobre a possibilidade da soma das partes não resultar em 100 \% por questões de \textbf{arredondamento}.
O percentual de alunos que consideraram a prova difícil ou muito difícil foi decrescente em função dos quartos de desempenho : 42,7 \% no primeiro quarto e 21,6 \% no quarto, grupo de melhor desempenho na prova.
Nos quartos de desempenho intermediários, a proporção de alunos que consideraram a prova difícil ou muito difícil foi de 37,2 \% no 2º quarto e 31,2 \% no 3º quarto.
As diferenças entre todos os quartos são estatisticamente significativas.
Para todos os quartos de desempenho, a alternativa modal para esta pergunta foi médio, com 49,6 \% e 62,7 \% dos respondentes nos quartos extremos, primeiro e quarto, respectivamente, valores crescentes com o desempenho. 4.1.2 Componente de Conhecimento Específico Ao responderem à Questão 2 – “ Qual o grau de dificuldade desta prova na parte de Componente Específico? ” – 36,9 \% do grupo de estudantes classificaram -na como difícil ou muito difícil.
Além disso, o Componente de Conhecimento Específico da prova foi considerado com grau de dificuldade médio por 55,6 \% dos alunos ( Gráfico 4.3, Gráfico 4.4, e, no Anexo II, a Tabela II.2 ).
A análise das respostas dos estudantes quanto ao grau de dificuldade do Componente de Conhecimento Específico da prova, agregado por Grande Região, mostra que a diferença entre a maior e as três menores proporções de alunos que a avaliaram como difícil ou muito difícil é estatisticamente significativa : a maior na região Sul ( 44,7 \% ) e as menores proporções nas regiões Sudeste ( 33,2 \% ), Nordeste ( 35,7 \% ) e Centro-Oeste ( 36,1 \% ).
O percentual de alunos que classificaram o grau de dificuldade como médio, no Componente de Conhecimento Específico, variou de 50,2 \% a 57,8 \%, para as regiões Sul e Nordeste, respectivamente.
Considerando -se a avaliação da dificuldade das questões do Componente de Conhecimento Específico da prova, de acordo com o desempenho dos estudantes, observa -se que há diferença estatisticamente significativa do resultado do quarto de melhor desempenho em relação ao terceiro quarto, e deste em relação ao segundo.
A proporção dos que classificaram a parte específica como difícil ou muito difícil variou de 29,4 \% ( 4º quarto ) a 42,9 \% ( 1º quarto ), valores decrescentes com o desempenho.
A alternativa modal para a Questão 2 foi o grau médio, com 49,4 \% do quarto inferior e 61,1 \% do superior optando por esta resposta, valores crescentes com o desempenho. 4.2 EXTENSÃO DA PROVA EM RELAÇÃO AO TEMPO TOTAL \textbf{Indagados} quanto à extensão da prova, em relação ao tempo total oferecido para a sua resolução ( Questão 3 ), os estudantes apontaram, com maior incidência, a alternativa que considerava a extensão adequada, para todas as agregações consideradas ( Gráfico 4.5, Gráfico 4.6, e, no Anexo II, a Tabela II.3 ).
O percentual de alunos que responderam ser a extensão da prova adequada foi de 37,6 \%.
Já 54,8 \% dos inscritos presentes consideraram que o exame foi longo ou muito longo e 7,6 \% o avaliaram como curto ou muito curto.
Entre as Grandes Regiões, a proporção daqueles que avaliaram a prova como longa ou muito longa em relação ao tempo total destinado à sua resolução variou pouco : de 46,6 \% na região Norte até 58,5 \% na região Sudeste.
A diferença entre as regiões Norte e Sudeste é estatisticamente significativa.
Considerando -se o desempenho dos alunos, nota -se ainda que 36,2 \% consideraram a extensão da prova adequada no quarto de desempenho inferior e 40,4 \% no de melhor desempenho ( quarto superior ), valores crescentes com o desempenho.
Nos quartos intermediários, esta proporção foi 35,5 \% no segundo quarto e 38,0 \% no terceiro.
No Gráfico 4.6, pode -se constatar para os que consideraram a prova longa ou muito longa em função dos quartos de desempenho que não há um padrão discernível.
Observa -se também que não há diferença estatisticamente significativa entre as proporções de escolha desta resposta entre os quartos. 4.3 COMPREENSÃO DOS ENUNCIADOS DAS QUESTÕES 4.3.1 Componente de Formação Geral Com relação aos enunciados das questões do Componente de Formação Geral ( Questão 4 ), as opiniões foram positivas, já que 71,3 \% dos alunos avaliados consideraram os enunciados de todas ou da maioria das questões claros e objetivos ( Gráfico 4.7, Gráfico 4.8, e, no Anexo II, a Tabela II.4 ).
Na análise regional, a \textbf{percentagem} de estudantes que avaliaram que todos ou a maioria dos enunciados das questões do Componente de Formação Geral estavam claros e objetivos variou de 69,1 \% na região Norte a 74,6 \% na região Sudeste, sendo esta diferença entre Norte e Sudeste estatisticamente significativa.
A análise das percepções dos estudantes sobre a clareza e objetividade dos enunciados permite afirmar que todos, ou a maioria dos enunciados de questões relativas ao Componente de Formação Geral, foram considerados claros e objetivos para a maior parte dos respondentes ( maior do que 69 \% em todas as regiões e maior ou igual do que 59 \% para todos os quartos de desempenho ).
Segundo o desempenho, observa -se que a proporção dos que emitiram esta opinião \textbf{cresce} conforme o desempenho aumenta, com diferenças estatisticamente significativas entre todos os quartos.
No quarto superior, a clareza e objetividade de todos ou da maioria dos enunciados das questões foi percebida por 82,2 \% dos alunos.
Já no quarto inferior esta proporção foi 59,0 \%. 4.3.2 Componente de Conhecimento Específico Com relação aos enunciados das questões do Componente de Conhecimento Específico da prova, para 70,3 \% dos estudantes avaliados da Área de Letras - Português, a clareza e a objetividade ( Questão 5 ) estavam presentes em todas ou na maioria das questões ( Gráfico 4.9, Gráfico 4.10, e no Anexo II, a Tabela II.5 ).
A maioria dos estudantes de todas as Grandes Regiões brasileiras considerou claros e objetivos todos ou a maioria dos enunciados das questões do Componente de Conhecimento Específico da prova, percentual sempre maior ou igual a 68,7 \%.
A diferença entre a maior e a menor proporção, 74,7 \% ( Centro-Oeste ) e 68,7 \% ( Nordeste ), é estatisticamente significativa.
A proporção de estudantes que consideraram os enunciados das questões claros e objetivos apresenta uma tendência crescente em relação ao aumento de desempenho : mais elevada no quarto superior ( 80,5 \% ) se comparada ao quarto inferior de desempenho ( 59,1 \% ).
As diferenças entre todos os quartos são estatisticamente significativas. 4.4 SUFICIÊNCIA DAS INFORMAÇÕES/INSTRUÇÕES FORNECIDAS Ao avaliarem as informações/instruções fornecidas para a resolução das questões ( Questão 6 ), 78,8 \% dos respondentes da Área de Letras - Português de todo o Brasil afirmaram que estas eram até excessivas ou suficientes em todas ou na maioria das questões ( Gráfico 4.11, Gráfico 4.12, e, no Anexo II, a Tabela II.6 ).
Quanto à distribuição de respondentes pelas Grandes Regiões observa -se que a proporção de estudantes que consideraram as informações/instruções fornecidas até excessivas ou suficientes em todas ou na maioria das questões foi sempre maior ou igual a 77,0 \%, chegando a 83,0 \% na região Sudeste.
A diferença entre o maior percentual, obtido na região Sudeste ( 83,0 \% ), e o menor, da região Nordeste ( 77,0 \% ) é estatisticamente significativa.
O percentual de participantes que avaliaram as informações/instruções como até excessivas ou suficientes em todas ou na maioria das questões foi mais elevado no quarto superior ( 88,8 \% ), percentual superior à média nacional ( 78,8 \% ).
Já no quarto inferior, a suficiência das informações/instruções declarada como até excessiva em todas ou na maioria das questões foi percebida por 68,1 \% dos respondentes.
Estas proporções são crescentes com os quartos de desempenho.
Levando -se em conta o desempenho dos participantes, notam -se diferenças estatisticamente significativas entre as opiniões de estudantes de todos os quartos, como mostra o Gráfico 4.12. 4.5 DIFICULDADE ENCONTRADA AO RESPONDER À PROVA Perguntados sobre as dificuldades com as quais se depararam ao responder à prova ( Questão 7 ), 15,8 \% dos estudantes apontaram o desconhecimento do conteúdo.
Para 46,7 \%, a forma diferente de abordagem do conteúdo foi indicada como dificuldade.
Já a falta de motivação para fazer a prova foi a dificuldade apontada por 18,8 \% dos respondentes.
Considerando -se todo o Brasil, 13,1 \% dos respondentes afirmaram que não tiveram qualquer tipo de dificuldade para responder à prova ( Tabela II.7 no Anexo II ).
Os Gráficos 4.13 e 4.14 apresentam os percentuais de estudantes que apontaram o desconhecimento do conteúdo como dificuldade percebida ao responder à prova.
Na análise por Grandes Regiões, o percentual de inscritos e presentes que apontaram o desconhecimento do conteúdo como dificuldade ao responder à prova não superou 16,8 \%.
Os percentuais variaram de 13,9 \% na região Sudeste a 16,8 \% na Norte, sendo que esta diferença não é estatisticamente significativa.
A forma diferente de abordagem do conteúdo foi a escolha modal dos estudantes, com percentuais que variaram de 41,8 \% ( região Sudeste ) a 51,5 \% ( Norte ).
O percentual de alunos que citaram a falta de motivação como dificuldade variou de 14,5 \% ( região Norte ) a 23,0 \% ( região Centro-Oeste ).
Os que declararam não ter qualquer dificuldade para responder à prova variaram de 10,2 \% na região Norte a 16,3 \% na Sudeste.
Com relação aos quartos de desempenho, o desconhecimento do conteúdo foi a opção escolhida por 15,0 \% dos estudantes do quarto superior e 15,4 \% do quarto inferior.
A alternativa modal para os alunos, quando agregados pelos quartos de desempenho, foi que a dificuldade encontrada foi causada pela forma diferente de abordagem do conteúdo : 47,3 \% no quarto inferior e 42,6 \% do quarto superior assim o responderam. 4.6 CONTEÚDOS DAS QUESTÕES OBJETIVAS DA PROVA Ao analisarem os conteúdos das questões objetivas da prova ( Questão 8 ), um percentual muito pequeno dos estudantes avaliados, apenas 8,8 \%, afirmou que não estudou ainda a maioria desses conteúdos ( Gráficos 4.15, Gráfico 4.16, e a Tabelas II.8 no Anexo II ).
A maioria ( 62,0 \% ) afirmou ter estudado e aprendido muitos ou todos os conteúdos.
Na análise por Grande Região, a proporção de respondentes que escolheram a opção não estudou ainda a maioria desses conteúdos foi pequena.
Observa -se que nas regiões Norte ( 11,9 \% ), Nordeste ( 9,2 \% ) e Sul ( 9,3 \% ), apesar de pequenas, as proporções foram maiores do que a média nacional ( 8,8 \% ).
Observa -se diferença estatisticamente significativa entre as regiões Norte e Nordeste, não ocorrendo o mesmo em relação à região Sul.
Em todas as regiões, a maioria dos presentes afirmou ter estudado e aprendido muitos ou todos os conteúdos, com proporções variando entre 56,3 \% na região Norte e 71,7 \% na Sudeste.
Considerando -se separadamente as opiniões de estudantes dos quatro quartos de desempenho, observa -se que, no quarto inferior, 15,6 \% ofereceram como resposta que não estudou ainda a maioria desses conteúdos, sendo 3,5 \% os do quarto superior com a mesma resposta, valores decrescentes com o desempenho.
As diferenças entre os alunos que optaram por este motivo de dificuldade nos diferentes quartos são estatisticamente significativas.
Tendo em conta o quarto superior, 79,7 \% dos alunos afirmaram ter estudado e aprendido muitos ou todos os conteúdos.
No outro extremo, no primeiro quarto, 44,7 \% optaram pelas mesmas categorias, valores crescentes com o desempenho. 4.7 TEMPO GASTO PARA CONCLUIR A PROVA Ao responderem sobre o tempo de conclusão da prova ( Questão 9 ), quase três quartos dos estudantes ( 70,0 \% ) afirmou ter gasto entre duas e quatro horas ( Gráfico 4.17, Gráfico 4.18 e, no Anexo II, a Tabela II.9 ).
Considerando -se as cinco Grandes Regiões brasileiras, o percentual dos que utilizaram entre duas e quatro horas para finalizar a prova nas regiões Centro-Oeste ( 62,3 \% ), Norte ( 68,2 \% ) e Sudeste ( 68,3 \% ), foi inferior ao percentual nacional.
Nas demais Grandes Regiões o percentual de alunos que \textbf{dispensaram} entre duas e quatro horas para concluir a prova, não ultrapassou a 73,6 \%, como mostra o Gráfico 4.17.
Uma vez sendo analisadas as alternativas escolhidas pelos estudantes que se situam nos diferentes quartos de desempenho, observa -se uma tendência crescente : uma maior proporção de participantes no quarto superior declarou ter gasto entre duas e quatro horas para concluir a prova quando comparadas com os dos quartos inferiores.
A diferença entre o terceiro ( 73,5 \% ) e o último quarto ( 75,4 \% ) não é estatisticamente significativa.

% matched lemmas: arredondamento, crescer, dispensar, indagar, numérico, ordenar, percentagem, questionário, recorte, reprodução, subconjunto
\end{textsample}
